% Sources, proof dependencies and board coverage: authoring/notes/unit03-editorial.md.
\section*{1. Bases and the number of independent vectors}

\notetarget{note:03:basis}A \textbf{basis} of a real linear space \(V\) is a finite list
\(v_1,\ldots,v_r\) that is linearly independent and spans \(V\).
The empty list is a basis of \(\{0\}\): its only linear combination is zero,
and it has no nonzero coefficient relation.
For a basis, every \(v\in V\) has exactly one expression
\[
v=c_1v_1+\cdots+c_rv_r.
\]
Existence is the spanning property. Subtracting two expressions gives
\(\sum_i(c_i-d_i)v_i=0\); independence gives \(c_i=d_i\) for all \(i\).
The standard columns \(e_1,\ldots,e_n\) form a basis of \(\R^n\), since
\(x=\sum_i x_ie_i\) and coordinate \(i\) of \(\sum_j c_je_j\) is \(c_i\).

\begin{example}\label{example:03:coordinates}\notetarget{note:03:coordinates}
Show that \(u=(1,1)^\top\) and \(v=(-1,2)^\top\) form a basis of \(\R^2\),
and express an arbitrary \((a,b)^\top\) in this basis.

\notetarget{note:03:coordinates-solution}The coefficient equations are
\[
x-y=a,\qquad x+2y=b.
\]
Subtracting the first from the second gives \(3y=b-a\), so
\[
\boxed{\begin{pmatrix}a\\b\end{pmatrix}
=\frac{2a+b}{3}\begin{pmatrix}1\\1\end{pmatrix}
+\frac{b-a}{3}\begin{pmatrix}-1\\2\end{pmatrix}.}
\]
These coefficients exist for every \(a,b\). For \(a=b=0\) both are zero,
which proves independence as well as spanning.
\end{example}

\begin{proposition}[Adjoining one vector]\label{prop:03:adjoin}\notetarget{note:03:adjoin}
If \(v_1,\ldots,v_r\) are independent, then
\[
v_1,\ldots,v_r,w\text{ are independent}
\quad\Longleftrightarrow\quad
w\notin\Span(v_1,\ldots,v_r).
\]
\end{proposition}
\begin{proof}\notetarget{note:03:adjoin-proof}
\begin{itemize}
\item If \(w=\sum_i a_iv_i\), then \(\sum_i a_iv_i-w=0\) has
coefficient \(-1\) on \(w\), so the enlarged list is dependent.
\item Suppose \(w\notin\Span(v_1,\ldots,v_r)\) and
\(\sum_i c_iv_i+dw=0\). If \(d\ne0\), then
\(w=-\sum_i(c_i/d)v_i\), a contradiction.
Thus \(d=0\); independence of the original list gives every \(c_i=0\).
\end{itemize}
\end{proof}

\begin{theorem}[Independent lists and spanning lists]\label{thm:03:bound}\notetarget{note:03:bound}
If \(v_1,\ldots,v_m\) span \(V\), every independent list in \(V\)
has at most \(m\) members.
\end{theorem}
\begin{proof}\notetarget{note:03:bound-proof}
\begin{itemize}
\item If \(m=0\), then \(V=\{0\}\), so every nonempty list is dependent.
\item Suppose \(m>0\) and \(w_1,\ldots,w_k\in V\), with \(k>m\). Write their coefficients
in columns of an \(m\times k\) matrix \(C=(c_{ij})\):
\[
w_j=\sum_{i=1}^m c_{ij}v_i\qquad(1\leqslant j\leqslant k).
\]
\item By Unit 1, Proposition~\ref{one-prop:01:excess}, \(Ct=0\) has a nonzero
solution \(t\in\R^k\). Regrouping the finite sum gives
\[
\sum_{j=1}^k t_jw_j
=\sum_{i=1}^m\left(\sum_{j=1}^k c_{ij}t_j\right)v_i=0.
\]
Thus the \(w_j\) are dependent.
\end{itemize}
\end{proof}

\begin{theorem}[Extracting and extending a basis]\label{thm:03:extract}\notetarget{note:03:extract}
Let \(V\) be spanned by a finite list.
\begin{enumerate}
\item A basis can be extracted from any finite spanning list.
\item Any independent list in \(V\) can be extended to a basis of \(V\).
\item Every subspace \(W\subseteq V\) has a finite basis.
\end{enumerate}
\end{theorem}
\begin{proof}\notetarget{note:03:extract-proof}
\begin{itemize}
\item Scan the spanning list from left to right, retaining a vector exactly
when it lies outside the span of those already retained.
Proposition~\ref{prop:03:adjoin} preserves independence at each step.
Every discarded vector belongs to the final span, so the retained list
still spans \(V\).
\item To retain specified independent vectors, place them first, followed
by a spanning list of \(V\), and apply the same procedure. None of the
specified vectors is discarded.
\item Starting with the empty list, choose a vector of \(W\) outside the
current span whenever one exists. The list stays independent.
If \(V\) is spanned by \(m\) vectors, Theorem~\ref{thm:03:bound}
forces the procedure to stop after at most \(m\) choices.
At that point the retained list spans \(W\). If \(W=\{0\}\), it stops
without choosing a vector.
\end{itemize}
\end{proof}

\section*{2. Dimension and adapted bases}

\notetarget{note:03:dimension}Any two bases of \(V\) have the same number of members:
if their lengths are \(r,s\), Theorem~\ref{thm:03:bound} gives both
\(r\leqslant s\) and \(s\leqslant r\).
This common number is the \textbf{dimension}, denoted \(\dim V\).
In particular \(\dim\{0\}=0\) and \(\dim\R^n=n\).

\notetarget{note:03:matrix-polynomial-bases}Let \(E_{ij}\) be the \(m\times n\) matrix with its only nonzero entry,
equal to 1, in position \((i,j)\). Then
\[
A=\sum_{i=1}^m\sum_{j=1}^n a_{ij}E_{ij}.
\]
Reading entry \((i,j)\) proves uniqueness, so the \(mn\) matrices \(E_{ij}\)
form a basis and \(\dim\R^{m\times n}=mn\).
In the polynomial coefficient space \(\mathcal P_d\) of Unit 2,
\(1,t,\ldots,t^d\) form a basis by the same coefficient argument;
thus \(\dim\mathcal P_d=d+1\).

\begin{proposition}[Dimension tests]\label{prop:03:dimension-tests}\notetarget{note:03:dimension-tests}
Let \(\dim V=n\) and \(W\subseteq V\) be a subspace.
\begin{enumerate}
\item Every independent list of \(n\) vectors in \(V\) is a basis.
\item Every spanning list of \(n\) vectors in \(V\) is a basis.
\item \(\dim W\leqslant n\), with equality exactly when \(W=V\).
\end{enumerate}
\end{proposition}
\begin{proof}\notetarget{note:03:dimension-tests-proof}
\begin{itemize}
\item Extend the independent list to a basis. Every basis has \(n\)
members, so no vector can have been added.
\item Extract a basis from the spanning list. Its length is \(n\), so
no vector can have been discarded.
\item A basis of \(W\) is independent in \(V\), giving the inequality.
At equality it is a basis of \(V\) by (a), and hence \(W=V\).
The reverse implication follows from equality of the spaces.
\end{itemize}
\end{proof}

\begin{example}\label{example:03:extend}\notetarget{note:03:extend-example}
Let \(V\) be the column span of
\[
A=\begin{pmatrix}1&2&3\\-2&3&8\\5&1&-3\\-3&-4&-5\end{pmatrix}
=[\,a_1\ a_2\ a_3\,].
\]
Find a basis of \(V\), its dimension, and an extension to a basis of \(\R^4\).

\notetarget{note:03:extend-solution}The relation \(a_3=2a_2-a_1\) permits deletion of \(a_3\).
If \(sa_1+ta_2=0\), the first two coordinates give
\[
s+2t=0,\qquad -2s+3t=0\quad\Longrightarrow\quad 7t=0, s=0.
\]
Thus \((a_1,a_2)\) is a basis of \(V\) and \(\dim V=2\).
For an extension, use \((a_1,a_2,e_3,e_4)\).
In a zero combination of these four columns the first two coordinates
again force the coefficients of \(a_1,a_2\) to be zero;
coordinates 3 and 4 then force the other two coefficients to be zero.
Four independent columns in \(\R^4\) form a basis.
\end{example}

\begin{theorem}[Dimension of a sum]\label{thm:03:sum}\notetarget{note:03:sum}
For subspaces \(U,W\) of a finite-dimensional real space,
\[
\boxed{\dim(U+W)=\dim U+\dim W-\dim(U\cap W).}
\]
\end{theorem}
\begin{proof}\notetarget{note:03:sum-proof}
\begin{itemize}
\item Choose a basis \(c_1,\ldots,c_k\) of \(U\cap W\), and extend it to bases
\[
(c_1,\ldots,c_k,u_1,\ldots,u_p)\text{ of }U,\qquad
(c_1,\ldots,c_k,w_1,\ldots,w_q)\text{ of }W.
\]
\item The combined list \((c_1,\ldots,c_k,u_1,\ldots,u_p,w_1,\ldots,w_q)\)
spans \(U+W\): add the two basis representations and collect the \(c_i\) coefficients.
\item For independence, suppose
\[
\sum_i a_ic_i+\sum_j b_ju_j+\sum_\ell d_\ell w_\ell=0.
\]
Then
\[
\sum_\ell d_\ell w_\ell=-\sum_i a_ic_i-\sum_j b_ju_j
\in U\cap W=\Span(c_1,\ldots,c_k).
\]
Write this vector as \(\sum_i e_ic_i\). The relation
\(\sum_\ell d_\ell w_\ell-\sum_i e_ic_i=0\) in the basis of \(W\)
gives every \(d_\ell=e_i=0\). The original relation now lies in the
basis of \(U\), giving every \(a_i=b_j=0\).
\item The combined basis has \(k+p+q=(k+p)+(k+q)-k\) members.
\end{itemize}
\end{proof}

\notetarget{note:03:intersection-bound}Since \(U+W\subseteq V\), this also gives
\[
\dim(U\cap W)\geqslant \dim U+\dim W-\dim V.
\]
In particular, if \(\dim U+\dim W>\dim V\), the intersection contains a
nonzero vector. For \(U\cap W=\{0\}\), the formula becomes
\(\dim(U\oplus W)=\dim U+\dim W\).

\begin{example}\label{example:03:matrix-spaces}\notetarget{note:03:matrix-spaces}
In \(\R^{n\times n}\), find bases and dimensions for the upper triangular
space \(T\), the symmetric space \(S\), and \(T\cap S\); determine \(T+S\).

\notetarget{note:03:matrix-spaces-solution}A basis of \(T\) is \(E_{ij}\), \(i\leqslant j\), since an upper
triangular matrix has the unique expression \(\sum_{i\leqslant j}a_{ij}E_{ij}\).
A symmetric matrix has the unique expression
\[
\sum_i a_{ii}E_{ii}+\sum_{i<j}a_{ij}(E_{ij}+E_{ji}).
\]
Hence a basis of \(S\) consists of \(E_{ii}\) and \(E_{ij}+E_{ji}\), \(i<j\).
There are \(n\) diagonal positions and \(n(n-1)/2\) unordered distinct
pairs: each of the \(n(n-1)\) ordered pairs belongs to one pair of reversals.
Thus \(\dim T=\dim S=n(n+1)/2\).
An upper triangular symmetric matrix is diagonal, so
\[
T\cap S=\Span(E_{11},\ldots,E_{nn}),\qquad \dim(T\cap S)=n.
\]
The sum formula gives \(\dim(T+S)=n^2\). Since \(T+S\subseteq\R^{n\times n}\),
Proposition~\ref{prop:03:dimension-tests} gives \(T+S=\R^{n\times n}\).
\end{example}

\begin{proposition}[Complements]\label{prop:03:complement}\notetarget{note:03:complement}
Every subspace \(W\subseteq V\) of a finite-dimensional space has a
subspace \(U\) such that \(V=W\oplus U\).
\end{proposition}
\begin{proof}\notetarget{note:03:complement-proof}
\begin{itemize}
\item Extend a basis \((w_1,\ldots,w_r)\) of \(W\) to a basis
\((w_1,\ldots,w_r,u_1,\ldots,u_s)\) of \(V\), and set
\(U=\Span(u_1,\ldots,u_s)\).
\item The extended basis spans \(V\), so \(W+U=V\).
\item If \(\sum_i a_iw_i=\sum_j b_ju_j\), independence of the full list gives
all \(a_i=b_j=0\). Thus \(W\cap U=\{0\}\), proving the direct sum.
\end{itemize}
\end{proof}

\begin{example}\label{example:03:complements}\notetarget{note:03:complement-example}
Let \(W=\Span((2,1)^\top)\) in \(\R^2\). Show that both
\(U=\Span((0,1)^\top)\) and \(U'=\Span((1,1)^\top)\) are complements of \(W\).

\notetarget{note:03:complement-solution}For every \((x,y)^\top\),
\[
\begin{pmatrix}x\\y\end{pmatrix}
=\frac{x}{2}\begin{pmatrix}2\\1\end{pmatrix}
+\left(y-\frac{x}{2}\right)\begin{pmatrix}0\\1\end{pmatrix}
=(x-y)\begin{pmatrix}2\\1\end{pmatrix}
+(2y-x)\begin{pmatrix}1\\1\end{pmatrix}.
\]
The equations \(a(2,1)^\top=b(0,1)^\top\) and
\(a(2,1)^\top=b(1,1)^\top\) each force \(a=b=0\), by their two coordinates.
Both sums are therefore direct. Since \((0,1)^\top\notin U'\),
the two complements differ.
\end{example}

\section*{3. Rank and rank factorization}

\notetarget{note:03:rank-definition}For \(A\in\R^{m\times n}\), let \(\mathcal C(A)\) be its column space,
\(N(A)=\{x:Ax=0\}\), and define its \textbf{row space}, represented by columns,
as \(\mathcal R(A)=\mathcal C(A^\top)\subseteq\R^n\).
The column rank is \(\dim\mathcal C(A)\); the row rank is \(\dim\mathcal R(A)\).

\begin{theorem}[Equality of the two ranks]\label{thm:03:row-column}\notetarget{note:03:row-column}
The column rank and row rank of every real matrix are equal.
Their common value is denoted \(\rank A\).
\end{theorem}
\begin{proof}\notetarget{note:03:row-column-proof}
\begin{itemize}
\item If \(A=0\), both spaces are \(\{0\}\), so both ranks are zero.
\item For \(A\ne0\), let \(c=\dim\mathcal C(A)\).
Extract a basis \(b_1,\ldots,b_c\) from
the columns of \(A\) and put \(B=[\,b_1\ \cdots\ b_c\,]\).
Express each column of \(A\) in this basis and collect the coefficient columns in \(F\):
\[
A=BF,\qquad B:m\times c,\quad F:c\times n.
\]
\item Row \(i\) of \(A\) is \(\sum_{j=1}^c b_{ij}\operatorname{row}_j(F)\).
Hence \(\mathcal R(A)\subseteq\mathcal R(F)\), and
\[
\dim\mathcal R(A)\leqslant c=\dim\mathcal C(A).
\]
\item Applying the same inequality to \(A^\top\) reverses it:
\[
\dim\mathcal C(A)=\dim\mathcal R(A^\top)
\leqslant\dim\mathcal C(A^\top)=\dim\mathcal R(A).
\]
Thus the two ranks are equal.
\end{itemize}
\end{proof}

\notetarget{note:03:rank-basic}Consequently,
\[
0\leqslant\rank A\leqslant\min(m,n),\qquad
\rank A^\top=\rank A,\qquad
\rank A=0\Longleftrightarrow A=0.
\]
The bound follows because the column and row spaces lie in \(\R^m\) and
\(\R^n\). Rank zero means every column belongs to \(\{0\}\).
A matrix has \textbf{full column rank} if its \(n\) columns are independent,
equivalently \(\rank A=n\); it has \textbf{full row rank} if \(\rank A=m\).
These equivalences follow by extracting bases from the respective lists.

\begin{proposition}[Full rank factors]\label{prop:03:factorization}\notetarget{note:03:factorization}
If \(\rank A=r>0\), there is a factorization
\[
A=BF,\qquad B\in\R^{m\times r},\quad F\in\R^{r\times n},
\qquad\rank B=\rank F=r.
\]
For this factorization,
\[
\mathcal C(A)=\mathcal C(B),\qquad
\mathcal R(A)=\mathcal R(F),\qquad N(A)=N(F).
\]
\end{proposition}
\begin{proof}\notetarget{note:03:factorization-proof}
\begin{itemize}
\item Use the extracted-column construction in Theorem~\ref{thm:03:row-column}.
The basis columns give \(\rank B=r\) and \(\mathcal C(A)=\mathcal C(B)\).
\item At the \(r\) selected column positions, the corresponding columns of \(F\)
are \(e_1,\ldots,e_r\), by uniqueness of basis coefficients.
Thus \(\rank F=r\).
\item Since \(\mathcal R(A)\subseteq\mathcal R(F)\)
and both have dimension \(r\), they are equal.
\item Finally \(BFx=0\) implies \(Fx=0\) by independence of the columns of \(B\);
the converse follows by multiplication by \(B\).
\end{itemize}
\end{proof}
\notetarget{note:03:zero-factorization}This is a \textbf{rank factorization}.
For \(r=0\), it is the product of an \(m\times0\) and
a \(0\times n\) matrix: each entry is an empty sum, hence zero.

\begin{proposition}[Invertible multiplication]\label{prop:03:rank-invariance}\notetarget{note:03:rank-invariance}
If \(P\) and \(Q\) are invertible of orders \(m\) and \(n\), then
\[
\rank(PAQ)=\rank A.
\]
\end{proposition}
\begin{proof}\notetarget{note:03:rank-invariance-proof}
\begin{itemize}
\item The rows of \(PA\) are combinations of rows of \(A\);
\(A=P^{-1}(PA)\) gives the reverse inclusion. Therefore
\(\mathcal R(PA)=\mathcal R(A)\) and \(\rank(PA)=\rank A\).
\item The columns of \(AQ\) are combinations of columns of \(A\);
\(A=(AQ)Q^{-1}\) gives the reverse inclusion. Therefore
\(\mathcal C(AQ)=\mathcal C(A)\) and \(\rank(AQ)=\rank A\).
\item Apply both equalities: \(\rank(PAQ)=\rank(AQ)=\rank A\).
\end{itemize}
\end{proof}

\section*{4. Computing bases and all solutions}

\begin{theorem}[The bases revealed by elimination]\label{thm:03:pivots}\notetarget{note:03:pivots}
Suppose row reduction gives an echelon matrix \(H=EA\), where \(E\)
is invertible, with nonzero rows \(1,\ldots,r\) and pivot columns
\(j_1<\cdots<j_r\).
\begin{enumerate}
\item The \(r\) nonzero rows of \(H\), transposed, are a basis of \(\mathcal R(A)\).
\item The original columns \(a_{j_1},\ldots,a_{j_r}\) are a basis of \(\mathcal C(A)\).
\item In particular, \(\rank A=r\).
\end{enumerate}
\end{theorem}
\begin{proof}\notetarget{note:03:pivots-proof}
\begin{itemize}
\item The nonzero rows span the row space of \(H\), which equals \(\mathcal R(A)\)
by the proof of Proposition~\ref{prop:03:rank-invariance}.
In a zero combination of these rows, inspecting pivot position \(j_1\)
forces the first coefficient to be zero. Inspecting \(j_2,j_3,\ldots\)
in order forces all remaining coefficients to be zero. Thus the row rank is \(r\).

\item For the pivot columns of \(H\), inspect row \(r\), then row \(r-1\),
and continue upward in a zero combination. Each pivot is nonzero, so all
coefficients vanish. There are \(r=\rank H\) independent pivot columns;
they form a basis of \(\mathcal C(H)\).
\item Multiplication by \(E\) preserves every coefficient relation among columns
(Unit 2, Proposition~\ref{two-prop:02:relations}).
Thus the original pivot columns are independent. If
\(h_j=\sum_i c_ih_{j_i}\), then
\[
E\left(a_j-\sum_i c_ia_{j_i}\right)=0
\quad\Longrightarrow\quad a_j=\sum_i c_ia_{j_i}.
\]
They also span \(\mathcal C(A)\), so they form a basis.
\end{itemize}
\end{proof}

\notetarget{note:03:reduced-echelon}Starting at the last pivot, divide its row by the pivot and clear all
entries above it; continue upward. In each row being used, earlier pivot
columns are zero by echelon form, and later pivot columns are zero from
the preceding clearing steps. Thus all other pivot columns are preserved.
The resulting matrix \(R\) has
a 1 in each pivot position and zeros everywhere else in a pivot column;
it is a \textbf{reduced echelon matrix}. Its zero rows may be omitted when
recording equations. This construction does not assert uniqueness.

\begin{theorem}[Null-space basis]\label{thm:03:nullity}\notetarget{note:03:nullity}
Let \(R\) be a reduced echelon matrix for \(A\), with pivots
\(j_1,\ldots,j_r\) and free column set \(F\).
For each \(k\in F\), set
\[
z_k=e_k-\sum_{i=1}^r R_{ik}e_{j_i}.
\]
Then \((z_k)_{k\in F}\) is a basis of \(N(A)\). In particular,
\[
\boxed{\dim N(A)=n-r,\qquad \rank A+\dim N(A)=n.}
\]
\end{theorem}
\begin{proof}\notetarget{note:03:nullity-proof}
\begin{itemize}
\item In row \(i\), the equation \(Rx=0\) reads
\[
x_{j_i}+\sum_{k\in F}R_{ik}x_k=0.
\]
Thus every solution has the expression \(x=\sum_{k\in F}x_kz_k\),
and conversely each such expression satisfies every row.
\item The free coordinate \(\ell\in F\) of \(z_k\) is 1 if \(\ell=k\) and 0
otherwise, so a zero combination of the \(z_k\) has all coefficients zero.
\item Row operations preserve solutions, hence this is a basis of \(N(A)\).
There are \(n-r\) free columns. When \(F\) is empty, the basis is empty
and \(N(A)=\{0\}\); for \(A=0\), it is the standard basis of \(\R^n\).
\end{itemize}
\end{proof}

\begin{example}\label{example:03:rank-computation}\notetarget{note:03:rank-example}
For
\[
A=\begin{pmatrix}1&5&6\\2&6&8\\7&1&8\end{pmatrix},
\]
find its rank, bases for its column, row, and null spaces, a rank factorization,
and all solutions of \(Ax=b\) for an arbitrary \(b=(b_1,b_2,b_3)^\top\).

\notetarget{note:03:rank-example-reduction}Apply \(R_2\leftarrow R_2-2R_1\), \(R_3\leftarrow R_3-7R_1\),
then \(R_3\leftarrow R_3-\tfrac{17}{2}R_2\):
\[
A\longrightarrow
\begin{pmatrix}1&5&6\\0&-4&-4\\0&-34&-34\end{pmatrix}
\longrightarrow
\begin{pmatrix}1&5&6\\0&-4&-4\\0&0&0\end{pmatrix}
\longrightarrow
R=\begin{pmatrix}1&0&1\\0&1&1\\0&0&0\end{pmatrix}.
\]
The last step divides row 2 by \(-4\) and subtracts five times that row from row 1.
There are two pivots, in columns 1 and 2. Therefore
\[
\rank A=2,\qquad
\mathcal C(A)=\Span\left(\begin{pmatrix}1\\2\\7\end{pmatrix},
\begin{pmatrix}5\\6\\1\end{pmatrix}\right),\qquad
\mathcal R(A)=\Span\left(\begin{pmatrix}1\\0\\1\end{pmatrix},
\begin{pmatrix}0\\1\\1\end{pmatrix}\right).
\]
Each displayed list is a basis. The free coordinate is \(x_3=t\), giving
\[
N(A)=\Span\left(\begin{pmatrix}-1\\-1\\1\end{pmatrix}\right),\qquad
A=\begin{pmatrix}1&5\\2&6\\7&1\end{pmatrix}
\begin{pmatrix}1&0&1\\0&1&1\end{pmatrix}.
\]
The product reproduces each original column, including \(a_3=a_1+a_2\).

\notetarget{note:03:rank-example-system}Applying the same row operations to the right-hand side gives
\[
[\,A\mid b\,]\longrightarrow
\left[\begin{array}{ccc|c}
1&0&1&(5b_2-6b_1)/4\\
0&1&1&(2b_1-b_2)/4\\
0&0&0&b_3+10b_1-17b_2/2
\end{array}\right].
\]
Thus the system is consistent exactly when
\(20b_1-17b_2+2b_3=0\), and then its full solution set is
\[
x=\begin{pmatrix}(5b_2-6b_1)/4\\(2b_1-b_2)/4\\0\end{pmatrix}
+t\begin{pmatrix}-1\\-1\\1\end{pmatrix},\qquad t\in\R.
\]
\end{example}

\begin{proposition}[Rank tests for solutions]\label{prop:03:consistent}\notetarget{note:03:consistent}
For \(A\in\R^{m\times n}\),
\begin{enumerate}
\item \(Ax=b\) is consistent exactly when \(\rank[\,A\ b\,]=\rank A\).
\item A consistent system has a unique solution exactly when \(\rank A=n\).
\item \(Ax=b\) is consistent for every \(b\in\R^m\) exactly when \(\rank A=m\).
\end{enumerate}
\end{proposition}
\begin{proof}\notetarget{note:03:consistent-proof}
\begin{itemize}
\item The inclusion \(\mathcal C(A)\subseteq\mathcal C([\,A\ b\,])\)
is equality exactly when \(b\in\mathcal C(A)\).
For nested subspaces, equality is equivalent to equality of dimensions.
\item By Unit 1, Theorem~\ref{one-thm:01:translate}, every nonempty solution
set is \(x_0+N(A)\). It is a singleton exactly when \(\dim N(A)=0\),
equivalently \(\rank A=n\) by Theorem~\ref{thm:03:nullity}.
\item Solvability for every \(b\) says \(\mathcal C(A)=\R^m\),
equivalently \(\dim\mathcal C(A)=m\).
\end{itemize}
\end{proof}
